%%%%%%%%%%%%%%%%%%%%%%%%%%%%%%%%%%%%%
%                                   %
% Compile with XeLaTeX and biber    %
%                                   %
% Questions or comments:            %
%                                   %
% joshua dot mcneill at uga dot edu %
%                                   %
%%%%%%%%%%%%%%%%%%%%%%%%%%%%%%%%%%%%%

\documentclass{beamer}
  % Read in standard preamble (cosmetic stuff)
  %%%%%%%%%%%%%%%%%%%%%%%%%%%%%%%%%%%%%%%%%%%%%%%%%%%%%%%%%%%%%%%%
% This is a standard preamble used in for all slide documents. %
% It basically contains cosmetic settings.                     %
%                                                              %
% Joshua McNeill                                               %
% joshua dot mcneill at uga dot edu                            %
%%%%%%%%%%%%%%%%%%%%%%%%%%%%%%%%%%%%%%%%%%%%%%%%%%%%%%%%%%%%%%%%

% Beamer settings
% \usetheme{Berkeley}
\usetheme{CambridgeUS}
% \usecolortheme{dove}
% \usecolortheme{rose}
\usecolortheme{seagull}
\usefonttheme{professionalfonts}
\usefonttheme{serif}
\setbeamertemplate{bibliography item}{}

% Packages and settings
\usepackage{fontspec}
  \setmainfont{Charis SIL}
\usepackage{hyperref}
  \hypersetup{colorlinks=true,
              allcolors=blue}
\usepackage{graphicx}
  \graphicspath{{../../figures/}}
\usepackage[normalem]{ulem}
\usepackage{enumerate}

% Document information
\author{M. McNeill}
\title[FREN2001]{Français 2001}
\institute{\url{joshua.mcneill@uga.edu}}
\date{}

%% Custom commands
% Lexical items
\newcommand{\lexi}[1]{\textit{#1}}
% Gloss
\newcommand{\gloss}[1]{`#1'}
\newcommand{\tinygloss}[1]{{\tiny`#1'}}
% Orthographic representations
\newcommand{\orth}[1]{$\langle$#1$\rangle$}
% Utterances (pragmatics)
\newcommand{\uttr}[1]{`#1'}
% Sentences (pragmatics)
\newcommand{\sent}[1]{\textit{#1}}
% Base dir for definitions
\newcommand{\defs}{../definitions}


  % Bibliography
  \usepackage[backend=biber,style=apa]{biblatex}
    \addbibresource{../references/References.bib}

  % Document information
  \subtitle[Welcome]{Welcome}

  %% Custom commands
  % Subsection/frame titles
  \newcommand{\subtwoone}{What is linguistics?}
  \newcommand{\subtwotwo}{What is language?}
  \newcommand{\subtwothree}{\citeauthor{hockett_origin_1960}'s (\citeyear{hockett_origin_1960}) design features of language}
  \newcommand{\subtwofour}{What you know about your language}
  \newcommand{\subtwofive}{What else?}
  \newcommand{\subtwosix}{How we study this stuff}
  \newcommand{\subtwoseven}{What about writing?}

\begin{document}
  % Read in the standard intro slides (title page and table of contents)
  %%%%%%%%%%%%%%%%%%%%%%%%%%%%%%%%%%%%%%%%%%%%%%%%%%%%%%%%%%%%%%%%
% This is a standard set of intro slides used in for all slide %
% documents. It basically contains the title page and table of %
% contents.                                                    %
%                                                              %
% Joshua McNeill                                               %
% joshua dot mcneill at uga dot edu                            %
%%%%%%%%%%%%%%%%%%%%%%%%%%%%%%%%%%%%%%%%%%%%%%%%%%%%%%%%%%%%%%%%

\begin{frame}
  \titlepage
  \tiny{Office: % Basically a variable for office hours location
Gilbert 121\\
        Office hours: % Basically a variable for office hours
 lundi, mercredi, vendredi 10:10--11:10
}
\end{frame}

\begin{frame}
  \tableofcontents[hideallsubsections]
\end{frame}

\AtBeginSection[]{
  \begin{frame}
    \tableofcontents[currentsection,
                     hideallsubsections]
  \end{frame}
}


  \section{Course Information}
    \begin{frame}{Description}
      \begin{block}{}
        The \emph{scientific} study of language, emphasizing such topics as the organization of grammar, language in space and time, and the relationship between the study of language and other disciplines.
      \end{block}
    \end{frame}

    \begin{frame}{Objectives}
      \begin{block}{}
        By the end of this course, you will be able to:
        \begin{itemize}
          \item Use linguistic terminology accurately
          \item Discuss basic linguistic facts, theories, and methodologies
          \item Analyze samples of written or spoken language from a variety of world languages
          \item Understand both language change and how language supports all learning and communication
        \end{itemize}
      \end{block}
    \end{frame}

    \begin{frame}{Materials}
      \begin{block}{}
        \fullcitebib{dawson_language_2016}
      \end{block}
    \end{frame}

    \begin{frame}{Grading}
      \begin{block}{}
        \begin{tabular}{r l}
          Participation & 7.5\%\\
          Homework      & 45\% (6x7.5\% + 1x2.5\%)\\
          Exams         & 45\% (6x7.5\%)
        \end{tabular}
      \end{block}
      \begin{block}{}
        \begin{itemize}
          \item Everything is 7.5\% of your grade
          \item There will be 6 exams (5 regular + final)
        \end{itemize}
      \end{block}
    \end{frame}

    \begin{frame}{Participation}
      \begin{block}{Attendance}
        \begin{itemize}
          \item \emph{Not} taken this semester
          \item You should attend only \emph{one} meeting per ``week''
          \item
          \begin{itemize}
            \item Group 1: First day of each ``week''
            \item Group 2: Second day of each ``week''
            \item Group 3: Third day of each ``week''
          \end{itemize}
        \end{itemize}
      \end{block}
      \begin{block}{eLC Discussions}
        \begin{itemize}
          \item For each section of the semester, post one question in the relevant eLC discussion forum
          \item Each question is worth 1.25\% of your final grade.
        \end{itemize}
      \end{block}
    \end{frame}

    \begin{frame}{Make-up Work}
      \begin{block}{}
        \emph{Not accepted}
      \end{block}
    \end{frame}

    \begin{frame}{Schedule}
      \begin{block}{}
        Before the beginning of in person meeting, you will be expected to have:
        \begin{itemize}
          \item Read the relevant section(s) of the book
          \item Watched the video lectures
          \item Thought about what you did understand and didn't understand
          \item Turned in any assignment(s) that are due
        \end{itemize}
      \end{block}
    \end{frame}

  \section{Introduction}
    \subsection{\subtwoone}
      \begin{frame}{\subtwoone}
        \begin{definition}<2->
          % Linguistics
The \emph{scientific} study of how language works

        \end{definition}
        \begin{block}<3->{A subfield of}
          \begin{itemize}
            \item Psychology and cognitive science when we're talking about the structure of a language
            \item Sociology and anthropology when we're talking about use, variation, and change
          \end{itemize}
        \end{block}
      \end{frame}

    \subsection{\subtwotwo}
      \begin{frame}{\subtwotwo}
        \begin{definition}<2->{}
          \begin{itemize}
            \item \alert{Language}: % Language
Human communication

            \item \alert{\emph{A} language}: % A language
A systematic set of words and rules that can be used for human communication \emph{and that is unique to each individual}

          \end{itemize}
        \end{definition}
        \begin{alertblock}<3->{}
          In this class, references to things like ``English'' are shorthand for ``the set of features that happen to be shared in the languages of some group of people''
        \end{alertblock}
      \end{frame}

    \subsection{\subtwothree}
      \begin{frame}{\subtwothree}
        \only<1-2>{
          \begin{block}<1-2>{Mode of communication}
            % Mode of communication
Languages can be communicated in different ways

            \begin{itemize}
              \item speech, gesture, writing
            \end{itemize}
          \end{block}
          \begin{block}<2>{Semanticity}
            % Semanticity
Language that is communicated has some meaning or function

            \begin{itemize}
              \item \lexi{pizza} means something, and you can assume that \lexi{alveolar} means something
            \end{itemize}
          \end{block}
        }
        \only<3-4>{
          \begin{block}<3-4>{Pragmatic function}
            % Pragmatic function
Language that is communicated has some purpose

            \begin{itemize}
              \item I didn't tell you about \lexi{alveolar} for nothing
            \end{itemize}
          \end{block}
          \begin{block}<4>{Interchangeability}
            % Interchangeability
Language can be transmitted and interpreted

            \begin{itemize}
              \item This is happening right now
            \end{itemize}
          \end{block}
        }
        \only<5-6>{
          \begin{block}<5-6>{Cultural transmission}
            % Cultural transmission
Features of languages can be passed from person to person through interaction

            \begin{itemize}
              \item You will (hopefully) know what \lexi{alveolar} means by the end of this course
            \end{itemize}
          \end{block}
          \begin{block}<6>{Arbitrariness} % Book introduces symbols and icons here
            % Arbitrariness
The connection between the forms of a language's words and their meanings is arbitrary

            \begin{itemize}
              \item A strawberry is called \lexi{strawberry} in English but \lexi{fraise} in French
            \end{itemize}
          \end{block}
        }
        \only<7-8>{
          \begin{block}<7-8>{Discreteness}
            % Discreteness
The elements of a language are made up of smaller elements

            \begin{itemize}
              \item Consonants and vowels make up words, words make up sentences
            \end{itemize}
          \end{block}
          \begin{block}<8>{Displacement}
            % Displacement
Language can be used to communicate about things that are not present

            \begin{itemize}
              \item \uttr{Last week, I went to New Orleans.}
            \end{itemize}
          \end{block}
        }
        \only<9>{
          \begin{block}{Productivity}
            % Productivity
The elements of a language can be combined and recombined in novel ways

            \begin{itemize}
              \item You've probably never heard the sentence: \uttr{Strawberry farmers throw pizza in New Orleans.}
            \end{itemize}
          \end{block}
        }
      \end{frame}

    \subsection{\subtwofour}
      \begin{frame}{\subtwofour}
        \only<1>{
          \begin{columns}
            \column{0.5\linewidth}
              \begin{block}{Competence vs performance \parencite{chomsky_syntactic_2002}}
                \begin{itemize}
                  \item \alert{Competence}: % Competence
What you know about your language

                  \item \alert{Performance}: % Performance
What you actually produce when you use your language

                \end{itemize}
              \end{block}
            \column{0.5\linewidth}
              \includegraphics[scale=0.125]{performance_cat.jpg}
          \end{columns}
        }
        \only<2>{
          \begin{block}{Aspects of your language}
            \begin{itemize}
              \item Phonetics and Phonology
              \item Morphology
              \item Syntax
              \item Semantics and Pragmatics
            \end{itemize}
          \end{block}
        }
        \only<3-4>{
          \begin{block}<3-4>{Phonetics}
            You know what is and is not a speech sound
            \begin{itemize}
              \item A dog's bark is not a speech sound
            \end{itemize}
          \end{block}
          \begin{block}<4>{Phonology}
            You know what combinations of speech sounds are possible
            \begin{itemize}
              \item /pt/ at the beginning of a word is not possible in your language
            \end{itemize}
          \end{block}
        }
        \only<5-6>{
          \begin{block}<5-6>{Morphology}
            You know where word boundaries are and the boundaries of the elements of words
            \begin{itemize}
              \item thedogisplayinginthebackyard
              % Might want to add
                % unbelievable has how many parts?
                % ungiraffelike makes sense? How?
            \end{itemize}
          \end{block}
          \begin{block}<6>{Syntax}
            You know what is and is not a grammatical (i.e., acceptable as possible) sentence
            \begin{itemize}
              \item \uttr{Colorless green ideas sleep furiously} vs
              \item \uttr{Furiously sleep ideas green colorless} \parencite{chomsky_syntactic_2002}
            \end{itemize}
          \end{block}
        }
        \only<7-8>{
          \begin{block}<7-8>{Semantics}
            You know the meanings of the words in your language
            \begin{itemize}
              \item What does \lexi{duck} mean?
            \end{itemize}
          \end{block}
          \begin{block}<8>{Pragmatics}
            You know what the purpose is behind an utterance you make
            \begin{itemize}
              \item What do you mean if you yell \lexi{duck}?
            \end{itemize}
          \end{block}
        }
        \only<9>{
          \begin{block}{}
            All the aspects of your language are stored in your \alert{mental lexicon} and your \alert{mental grammar}
            \begin{itemize}
              \item Mental lexicon: % Mental lexicon
All the words you know, their meanings, functions, pronunciations, and relations to other words

              \item Mental grammar: % Mental grammar
The set of rules that you use to combine sounds, parts of words, and words

            \end{itemize}
          \end{block}
        }
      \end{frame}

    \subsection{\subtwofive}
      \begin{frame}{\subtwofive}
        \begin{block}{Psycholinguistics}
          % Psycholinguistics
The study of how the mind processes, learns, stores, and produces language

          \begin{itemize}
            \item This is the field that tests theories about what you know about your language
          \end{itemize}
        \end{block}
        \begin{block}<2->{Language variation and change}
          % Language variation and change
How the languages of individuals within and between communities differ and how the languages of individuals of different time periods differ

          \begin{itemize}
            \item Studied in sociolinguistics
          \end{itemize}
        \end{block}
      \end{frame}

      \begin{frame}{\subtwofive}
        \begin{block}{Computational linguistics}
          % Computational linguistics
The study of how to engineer computers to process and produce language

        \end{block}
      \end{frame}

    \subsection{\subtwosix}
      \begin{frame}{\subtwosix}
        \only<1-3>{
          \begin{block}{Mental grammar}
            % Mental grammar
The set of rules that you use to combine sounds, parts of words, and words

          \end{block}
          \begin{block}<2->{Descriptive grammar}
            % Descriptive grammar
The set of rules that you deduce from observing a speaker's linguistic performance

          \end{block}
          \begin{block}<3->{Prescriptive grammar}
            % Prescriptive grammar
The set of rules that are prescribed for some context

          \end{block}
        }
        \only<4>{
          \begin{example}
            A speaker says, \uttr{Where do you come \emph{from}?}, which contains a stranded preposition
            \begin{itemize}
              \item Mental grammar: There is definitely some rule about this in this speaker's head
              \item Descriptive grammar: This speaker's language seems to allow stranded prepositions
              \item Prescriptive grammar: Stranded prepositions are inappropriate when writing essays
            \end{itemize}
          \end{example}
        }
      \end{frame}

    \subsection{\subtwoseven}
      \begin{frame}{\subtwoseven}
        \begin{alertblock}{}
          Writing is not \emph{the} language nor a more perfect form of language nor primary over speech
        \end{alertblock}
      \end{frame}

  \begin{frame}{References}
    \printbibliography
  \end{frame}
\end{document}
